% Options for packages loaded elsewhere
\PassOptionsToPackage{unicode}{hyperref}
\PassOptionsToPackage{hyphens}{url}
\PassOptionsToPackage{dvipsnames,svgnames,x11names}{xcolor}
\documentclass[
  10pt,
]{article}
\usepackage{xcolor}
\usepackage[margin=2.4cm]{geometry}
\usepackage{amsmath,amssymb}
\setcounter{secnumdepth}{-\maxdimen} % remove section numbering
\usepackage{iftex}
\ifPDFTeX
  \usepackage[T1]{fontenc}
  \usepackage[utf8]{inputenc}
  \usepackage{textcomp} % provide euro and other symbols
\else % if luatex or xetex
  \usepackage{unicode-math} % this also loads fontspec
  \defaultfontfeatures{Scale=MatchLowercase}
  \defaultfontfeatures[\rmfamily]{Ligatures=TeX,Scale=1}
\fi
\usepackage{lmodern}
\ifPDFTeX\else
  % xetex/luatex font selection
  \setmainfont[]{Libertinus Serif}
  \setmonofont[Scale=0.85]{Latin Modern Mono}
  \setmathfont[]{Latin Modern Math}
\fi
% Use upquote if available, for straight quotes in verbatim environments
\IfFileExists{upquote.sty}{\usepackage{upquote}}{}
\IfFileExists{microtype.sty}{% use microtype if available
  \usepackage[]{microtype}
  \UseMicrotypeSet[protrusion]{basicmath} % disable protrusion for tt fonts
}{}
\makeatletter
\@ifundefined{KOMAClassName}{% if non-KOMA class
  \IfFileExists{parskip.sty}{%
    \usepackage{parskip}
  }{% else
    \setlength{\parindent}{0pt}
    \setlength{\parskip}{6pt plus 2pt minus 1pt}}
}{% if KOMA class
  \KOMAoptions{parskip=half}}
\makeatother
\usepackage{longtable,booktabs,array}
\newcounter{none} % for unnumbered tables
\usepackage{calc} % for calculating minipage widths
% Correct order of tables after \paragraph or \subparagraph
\usepackage{etoolbox}
\makeatletter
\patchcmd\longtable{\par}{\if@noskipsec\mbox{}\fi\par}{}{}
\makeatother
% Allow footnotes in longtable head/foot
\IfFileExists{footnotehyper.sty}{\usepackage{footnotehyper}}{\usepackage{footnote}}
\makesavenoteenv{longtable}
\setlength{\emergencystretch}{3em} % prevent overfull lines
\providecommand{\tightlist}{%
  \setlength{\itemsep}{0pt}\setlength{\parskip}{0pt}}

\providecommand{\E}{\mathbb{E}}
\providecommand{\N}{\mathbb{N}}
\providecommand{\Z}{\mathbb{Z}}
\providecommand{\R}{\mathbb{R}}
\providecommand{\eps}{\varepsilon}
\providecommand{\ceil}[1]{\left\lceil #1\right\rceil}
\providecommand{\floor}[1]{\left\lfloor #1\right\rfloor}
\AtBeginDocument{\let\setminus\smallsetminus}
\usepackage[most]{tcolorbox}
\newtcolorbox{warning}{colback=red!5,colframe=red!65!black,boxrule=0.8pt,arc=2pt,left=6pt,right=6pt,top=5pt,bottom=5pt}
\usepackage{etoolbox}
\AtBeginEnvironment{longtable}{\small}
\setlength{\emergencystretch}{3em}
\usepackage{fancyhdr}
\pagestyle{fancy}
\fancyhf{}
\fancyhead[L]{\small\bfseries UNREFEREED GPT-6 ASTRA OUTPUT}
\fancyhead[R]{\small\thepage}
\renewcommand{\headrulewidth}{0.4pt}
\usepackage{bookmark}
\IfFileExists{xurl.sty}{\usepackage{xurl}}{} % add URL line breaks if available
\urlstyle{same}
\hypersetup{
  pdftitle={Covering powers of cycles with equivalence subgraphs},
  pdfauthor={Writeup: gpt-6-astra (single model pass)},
  colorlinks=true,
  linkcolor={blue!50!black},
  filecolor={Maroon},
  citecolor={Blue},
  urlcolor={blue!50!black},
  pdfcreator={LaTeX via pandoc}}

\title{Covering powers of cycles with equivalence subgraphs}
\usepackage{etoolbox}
\makeatletter
\providecommand{\subtitle}[1]{% add subtitle to \maketitle
  \apptocmd{\@title}{\par {\large #1 \par}}{}{}
}
\makeatother
\subtitle{Unrefereed candidate disproof by GPT-6 Astra}
\author{Writeup: \texttt{gpt-6-astra} (single model pass)}
\date{Catalog id
\texttt{covering\_powers\_of\_cycles\_with\_equivalence\_subgraphs} ---
generated 2026-09-16}

\begin{document}
\maketitle

\begin{warning}

\textbf{UNREFEREED MODEL OUTPUT.} This document was selected solely
because GPT-6 Astra labelled its own result
\texttt{would\_publish:\ true} and returned \texttt{proved} or
\texttt{disproved}. It has not passed the adversarial LLM referee used
for the earlier Sol campaign, has not been checked by a human
mathematician, and has not been checked for novelty. Treat every
mathematical and bibliographic claim below as unverified.

\end{warning}

\section{Summary and provenance}\label{summary-and-provenance}

{\def\LTcaptype{none} % do not increment counter
\begin{longtable}[]{@{}
  >{\raggedright\arraybackslash}p{(\linewidth - 2\tabcolsep) * \real{0.5000}}
  >{\raggedright\arraybackslash}p{(\linewidth - 2\tabcolsep) * \real{0.5000}}@{}}
\toprule\noalign{}
\endhead
\bottomrule\noalign{}
\endlastfoot
Catalog id &
\texttt{covering\_powers\_of\_cycles\_with\_equivalence\_subgraphs} \\
Catalog entry &
\href{https://graph-theory-ai.github.io/graph-conjectures/op/covering_powers_of_cycles_with_equivalence_subgraphs/}{Covering
powers of cycles with equivalence subgraphs} \\
Source corpus & opg \\
Campaign leg & OpenProblemGarden first pass (\texttt{attacks\_opg}) \\
Source paper / entry &
http://www.openproblemgarden.org/op/covering\_powers\_of\_cycles\_with\_equivalence\_subgraphs \\
Model verdict & \textbf{disproved} (confidence: high) \\
Model's one-line claim & Noncomplete powers of cycles have equivalence
covering number Θ(log(k+1)), contradicting the proposed Ω(k) lower
bound. \\
Model & \texttt{gpt-6-astra}, reasoning effort \texttt{max},
\texttt{mode=pro}, flex service tier \\
Original artifact &
\texttt{attacks\_opg/covering\_powers\_of\_cycles\_with\_equivalence\_subgraphs/output.md} \\
Independent review & \textbf{None} \\
Model's caveats & Bibliographic novelty is unverified; an exact formula
is obtained when k+1 divides n, but not for every n. \\
\end{longtable}
}

\section{Problem statement}\label{problem-statement}

Conjecture Given \(k\) and \(n\) , the graph \(C_{n}^k\) has equivalence
covering number \(\Omega(k)\) .

\subsection{Catalog context}\label{catalog-context}

Given a graph \(G\) , a subgraph \(H\) of \(G\) is an equivalence
subgraph of \(G\) if \(H\) a disjoint union of cliques. The quivalence
covering number of \(G\) , denoted \(eq(G)\) , is the least number of
equivalence subgraphs needed to cover the edges of \(G\) . This problem
has been studied by various people since the 80s {[}A{]}. For line
graphs, the equivalence covering number is known to within a constant
factor {[}EGK{]}. It is therefore tempting to examine the situation for
quasi-line graphs and claw-free graphs. Powers of cycles are perhaps the
simplest interesting class of claw-free graphs that are not necessarily
line graphs. However, even for \(n\) very large compared to \(k\) , no
upper bound is known beyond trivial linear bounds of order \(\Theta(k)\)
. Furthermore, it is not even certain that a nontrivial lower bound
(i.e.~going to infinity as \(k\) goes to infinity) is known. It is
possible that this can be related somehow to a known result, but for now
it seems at least superficially that this problem is wide open.

\section{Astra writeup}\label{astra-writeup}

\emph{The text below is the model output verbatim, apart from moving its
machine-readable verdict block into the summary above and shifting
Markdown heading levels for this document. Mechanical TeX defects and
equation tags were normalized where needed for compilation.}

\subsubsection{Result}\label{result}

The linear conjecture is false under the definition in the question. The
following bounds are uniform in \(n\), so the counterexample does not
depend on taking \(n\) small relative to \(k\).

\textbf{Theorem.} Let \(k\ge 1\), put \[
s=k+1,\qquad r=\left\lceil\log_2 s\right\rceil.
\] For every \(n\ge 2k+2\), \[
\boxed{\quad r+1\le \operatorname{eq}(C_n^k)\le 2r+1.\quad}
\] Moreover, if \(s\mid n\), then \[
\boxed{\quad \operatorname{eq}(C_n^k)=r+1.\quad}
\]

The remaining regime \(n\le 2k+1\), with \(n\ge3\), consists of complete
graphs and has equivalence covering number \(1\).

In particular, for \[
k=2^q-1,\qquad n=m2^q,\qquad m\ge2,
\] the exact value is \[
\operatorname{eq}(C_n^k)=q+1.
\] Here \(m\) can be arbitrarily large. For example, taking \(m=2^q\)
gives \(n/k\to\infty\), while \[
\frac{\operatorname{eq}(C_n^k)}{k}
=\frac{q+1}{2^q-1}\longrightarrow0.
\]

All equivalence subgraphs below may be made spanning by adding isolated
vertices. Edges of the ambient graph between different chosen cliques
are simply omitted from that equivalence subgraph.

\subsubsection{1. Binary splitting of a
boundary}\label{binary-splitting-of-a-boundary}

Write each integer \(x\in\{0,\dots,s-1\}\) using \(r\) binary digits,
most significant first: \[
x=x_1x_2\cdots x_r.
\] For a position \(h\in\{1,\dots,r\}\), a binary prefix \(p\) of length
\(h-1\), and \(\varepsilon\in\{0,1\}\), define \[
I_h^\varepsilon(p)
=
\{x\in\{0,\dots,s-1\}:
(x_1,\dots,x_{h-1})=p,\ x_h=\varepsilon\}.
\]

Two elementary properties will be used:

\begin{enumerate}
\def\labelenumi{\arabic{enumi}.}
\tightlist
\item
  If \(x\in I_h^1(p)\) and \(y\in I_h^0(p)\), then \(x>y\).
\item
  Whenever \(x>y\), their first differing binary digit gives an \(h,p\)
  such that \[
  x\in I_h^1(p),\qquad y\in I_h^0(p).
  \]
\end{enumerate}

Thus the relation \(x>y\) is covered by these binary comparisons, using
\(r\) levels.

\subsubsection{\texorpdfstring{2. An \((r+1)\)-cover when
\(s\mid n\)}{2. An (r+1)-cover when s\textbackslash mid n}}\label{an-r1-cover-when-smid-n}

Suppose \(n=ms\), where \(m\ge2\). Partition the cycle into consecutive
blocks \[
B_i=\{v_{i,0},\dots,v_{i,s-1}\},
\qquad i\in\mathbb Z/m\mathbb Z.
\] The cyclic order runs through the local indices \(0,\dots,s-1\) in
each block.

Let \(H_0\) be the disjoint union of the cliques on the blocks \(B_i\).
These are valid cliques because each block consists of \(s=k+1\)
consecutive vertices.

For each binary position \(h\), and for every block \(i\) and prefix
\(p\), form \[
Q_{i,h,p}
=
\{v_{i,x}:x\in I_h^1(p)\}
\ \cup\
\{v_{i+1,y}:y\in I_h^0(p)\}.
\] Let \(H_h\) be the union of the complete graphs on these sets,
omitting empty sets.

\paragraph{\texorpdfstring{Each \(Q_{i,h,p}\) is a
clique}{Each Q\_\{i,h,p\} is a clique}}\label{each-q_ihp-is-a-clique}

Pairs in the same block are adjacent. For a cross-block pair, its local
indices satisfy \(x>y\). The clockwise distance from \(v_{i,x}\) to
\(v_{i+1,y}\) is \[
s+y-x\le s-1=k.
\] Hence every cross-block pair is also adjacent.

\paragraph{\texorpdfstring{The cliques in \(H_h\) are
vertex-disjoint}{The cliques in H\_h are vertex-disjoint}}\label{the-cliques-in-h_h-are-vertex-disjoint}

Fix \(h\). A vertex \(v_{i,x}\) has one binary prefix \(p\) of length
\(h-1\).

\begin{itemize}
\tightlist
\item
  If \(x_h=1\), it belongs to \(Q_{i,h,p}\).
\item
  If \(x_h=0\), it belongs to \(Q_{i-1,h,p}\).
\end{itemize}

It belongs to no other such set. Consequently, \(H_h\) is an equivalence
subgraph.

This is the key compatibility: at a given binary level, a vertex
participates across its right boundary or its left boundary, but never
both.

\paragraph{All edges are covered}\label{all-edges-are-covered}

An edge of \(C_n^{s-1}\) has a cyclic arc of length at most \(s-1\)
joining its endpoints. Such an arc cannot cross two block boundaries.

If its endpoints lie in one block, \(H_0\) covers it. Otherwise, orient
its short arc from \(B_i\) to \(B_{i+1}\), and write the endpoints as
\(v_{i,x},v_{i+1,y}\). Adjacency gives \[
s+y-x\le s-1,
\] so \(x>y\). At their first differing binary digit, the two endpoints
belong to the same \(Q_{i,h,p}\).

Therefore \[
\operatorname{eq}(C_{ms}^{s-1})\le r+1.
\]

Some of these cliques have gaps in the cyclic order. That is allowed:
they are cliques, though not necessarily consecutive or maximal cliques.

\subsubsection{\texorpdfstring{3. A uniform upper bound for arbitrary
\(n\)}{3. A uniform upper bound for arbitrary n}}\label{a-uniform-upper-bound-for-arbitrary-n}

We first record the linear version of the preceding construction.

\paragraph{Path-power cover}\label{path-power-cover}

For every \(N\), the graph on \(0,\dots,N-1\) with edges \[
uv\in E\quad\Longleftrightarrow\quad 0<|u-v|\le s-1
\] has an equivalence cover of size at most \(r+1\).

Indeed, partition the linear order into blocks of size \(s\), with the
last block possibly shorter. Use the block-clique layer \(H_0\), and use
the same binary construction across every consecutive pair of blocks. At
each binary level, the cliques remain vertex-disjoint by the same
left/right membership argument.

Every edge lies within one block or crosses one boundary. Across a
boundary, the local indices again satisfy \(x>y\), so the first
differing binary digit covers the edge. Missing vertices in the final
block cause no difficulty.

\paragraph{Closing the cycle}\label{closing-the-cycle}

Now label the vertices of \(C_n^{s-1}\) by \(0,\dots,n-1\), where
\(n\ge2s\).

The non-wrapping edges, namely those with ordinary index difference at
most \(s-1\), are covered by the preceding path-power construction using
\(r+1\) layers.

Every remaining edge wraps across the cut between \(n-1\) and \(0\).
Consider the disjoint sets \[
A=\{n-s,\dots,n-1\},\qquad
B=\{0,\dots,s-1\}.
\] Both are cliques. Write their vertices as \[
a_x=n-s+x,\qquad b_y=y,
\qquad 0\le x,y<s.
\] A wrap-around edge satisfies \[
n-a_x+b_y=s+y-x\le s-1,
\] equivalently \(x>y\). Conversely, every pair with \(x>y\) is an edge
across this boundary.

For each \(h,p\), take the clique \[
\{a_x:x\in I_h^1(p)\}
\ \cup\
\{b_y:y\in I_h^0(p)\}.
\] For fixed \(h\), these cliques are vertex-disjoint, and the \(r\)
binary levels cover all wrap-around edges.

Combining the two covers gives \[
\operatorname{eq}(C_n^{s-1})
\le (r+1)+r=2r+1.
\]

\subsubsection{4. A logarithmic lower bound by matrix
rank}\label{a-logarithmic-lower-bound-by-matrix-rank}

The following elementary rank argument proves the lower bound, including
the exact value in the divisible case.

Let \[
t=\operatorname{eq}(C_n^{s-1}),\qquad n\ge2s.
\] Since \(s\ge2\), this graph has edges and \(t\ge1\).

For each layer \(\ell\in\{1,\dots,t\}\), assign distinct real labels to
its clique components, and let \(c_\ell(v)\) be the label of the
component containing \(v\). Then, for distinct vertices, \[
uv\in E(C_n^{s-1})
\quad\Longleftrightarrow\quad
c_\ell(u)=c_\ell(v)\text{ for some }\ell.
\]

Define an \(n\times n\) real matrix \[
M_{uv}
=
\prod_{\ell=1}^{t}\bigl(c_\ell(u)-c_\ell(v)\bigr).
\] Thus \(M_{uv}=0\) exactly when \(u=v\) or \(uv\) is an edge.

\paragraph{Rank upper bound}\label{rank-upper-bound}

For \(S\subseteq[t]\), let \[
a_S(v)=\prod_{\ell\in S}c_\ell(v),
\] using the empty-product convention. Expanding the product gives \[
M
=
\sum_{S\subseteq[t]}
(-1)^{t-|S|}\,a_Sa_{[t]\setminus S}^{\,T}.
\] Each summand has rank at most one, so \[
\operatorname{rank}M\le2^t.
\]

\paragraph{\texorpdfstring{A nonsingular minor of order
\(2s\)}{A nonsingular minor of order 2s}}\label{a-nonsingular-minor-of-order-2s}

Use cyclic vertex labels modulo \(n\), and consider the ordered
submatrix \[
N_{ij}=M_{\,i,\,(s+j)\bmod n},
\qquad 0\le i,j<2s.
\] Its row indices are distinct and its column indices are distinct,
since \(2s\le n\).

If \(i>j\), then \[
1-s\le s+j-i\le s-1.
\] Therefore the corresponding vertices are equal or have cyclic
distance at most \(s-1\), and \[
N_{ij}=0.
\]

On the diagonal, the two vertices have cyclic distance exactly \(s\),
because \(n\ge2s\). They are distinct and nonadjacent, so \[
N_{ii}\ne0.
\]

Hence \(N\) is upper triangular with nonzero diagonal. It follows that
\[
2s=\operatorname{rank}N
\le\operatorname{rank}M
\le2^t.
\] Consequently, \[
t\ge\left\lceil\log_2(2s)\right\rceil
=1+\left\lceil\log_2s\right\rceil
=r+1.
\]

Together with the construction in Section 2, this proves \[
\operatorname{eq}(C_{ms}^{s-1})
=1+\left\lceil\log_2s\right\rceil
\qquad(m\ge2).
\]

\subsubsection{Conclusion and scope}\label{conclusion-and-scope}

The bounds establish \[
\operatorname{eq}(C_n^k)=\Theta(\log(k+1))
\] uniformly over all positive \(k\) and all \(n\ge2k+2\). Thus the
proposed \(\Omega(k)\) lower bound fails even when \(n\) is arbitrarily
large relative to \(k\).

No unproved conjecture or external theorem is used. The exact value for
every nondivisible \(n\) is not determined here. I have also not
verified whether this binary construction and rank argument already
appear in the literature; the conclusion is a self-contained disproof of
the stated conjecture, not a claim of historical novelty.

\end{document}
